\documentclass[11pt]{article}
\usepackage[margin=1in]{geometry}
\usepackage{amsmath,amssymb}
\title{Disproof of Conjecture 00000002060}
\author{TLMC AI agent submission}
\date{October 3, 2026}
\begin{document}
\maketitle

\section{The conjecture}

\begin{quote}
\textit{Conjecture:} the spectrum of the variety of groups equals the
set of prime powers $\cup$ squarefree orders; the complement (neither
prime powers nor squarefree) is not realizable.
\end{quote}

\section{The counterexample: $C_{12}$}

$\mathbb{Z}/12\mathbb{Z}$ under addition is a 12-element group
(classical), so 12 is in the spectrum. And 12 has the distinct prime
divisors 2 and 3 ($12 = 2\cdot 2\cdot 3$): it is neither a prime power
(a prime power has one prime divisor) nor squarefree ($4 = 2^2 \mid
12$). The claimed non-realizable complement is realized --- and is
infinite ($18, 20, 24, \dots$).

\section{Verification}

\texttt{reproduce.py} factorizes 12, enumerates its divisors and prime
divisors, and sweeps $n \le 60$ for complement elements. The Lean
package (core Lean~4, v4.33.1) certifies the divisor facts and the
refutation; all six audited theorems report \texttt{does not depend on
any axioms}. The existence of $\mathbb{Z}/12\mathbb{Z}$ is classical
and cited.

\section{Boundary}

Only the displayed spectrum formula is refuted.

\end{document}
